\section{An actual normal alternating core in the Heisenberg label action}
\label{sec:heisenberg-normal-alternating-core}

\paragraph{Author freeze and origin.}
D-TD37, reserved question Q2es, is a derivation by
\texttt{/root/imprimitive}.  Root proposed combining the actual
central-product branch with support-sensitive conjugate generation;
the present extension treats every isotypic orbit of an actual normal
quasisimple core, including compatible odd operators and permutations
of absolute constituents.  The inputs D-TD27, D-TD31 and D-TD32
have separate whole-proof reviews, and HD-FT27 is an accepted
actual-source obstruction.  Complete author proof, source-interface,
checker and first-output readings pass.  Independent and root whole-proof
reviews of this new section also pass.  No abstract quotient is
asserted to lift to a subgroup.

\subsection*{D-TD37. A rank comparison for every compatible label operator}

Let \(\mathcal Q\) be a nonzero finite-dimensional vector space over
\(\mathbf F_p\), put \(d=\dim_{\mathbf F_p}\mathcal Q\), and let
\(H\le\operatorname{GL}_{\mathbf F_p}(\mathcal Q)\) be its actual
faithful finite image, irreducible over this original field.  Assume
there is an actual normal subgroup and a compatible actual quotient
\[
 E\triangleleft H,\qquad E=E',\qquad
 E/Z(E)\simeq A_n,\qquad
 \psi:H\twoheadrightarrow T,\quad T\in\{A_n,S_n\},\quad n\ge5,
 \qquad \psi(E)=A_n,\quad\ker(\psi|_E)=Z(E).
 \tag{ANE.1}
\]
In particular the restriction of the given matrix representation to E
is faithful.  The natural action on n letters is the one fixed by
\(\psi\); it is not selected independently for each operator.
For \(h\in H\), let \(k(h)\) be the number of letters moved by
\(\psi(h)\), and define the geometric projective rank by
\[
 \nu_{\rm alg}(h)
   =d-\max_{\lambda\in\overline{\mathbf F}_p}
          \dim\ker(h-\lambda I\mid\mathcal Q\otimes\overline{\mathbf F}_p).
 \tag{ANE.2}
\]

\begin{theorem}[D-TD37]
Under (ANE.1), every h satisfies
\[
 \operatorname{rank}_{\mathbf F_p}(1-h)
       \ge\nu_{\rm alg}(h)\ge\frac{d\,k(h)}{24n},
       \qquad d>\frac n5.
 \tag{ANE.3}
\]
No absolute irreducibility, primitivity or tensor indecomposability
of the defining H-module is required.  No semisimplicity of an
individual h is required either.

Moreover, suppose this is the actual defining action of a stable
whole projective Heisenberg factor in a faithful finite complex
coefficient representation of a connected sphere cover.  Precisely,
suppose
\[
 \begin{gathered}
 C\longrightarrow\mathbf P^1\text{ is a connected }J\text{-Galois cover},
 \qquad V=W\otimes R,\qquad e=\dim V,\\
 P=\operatorname{im}(J\to\operatorname{PGL}(W)),\qquad
 Q\triangleleft P,\quad Q\simeq\mathcal Q,\qquad
 P/Q=H\le\operatorname{Sp}(\mathcal Q),\qquad
 a=\dim W=p^{d/2},
 \end{gathered}
 \tag{ANE.4}
\]
where the projective Q-action is irreducible with nondegenerate
scalar commutator form, and R is the entire external partner.
If \(n\ge10^6\), the actual coefficient inertia cannot have
moving-rank sum \(2e+6\).  There is no required lower bound on
the number of projected nonidentity marks and no required support
floor on R.
\end{theorem}

\begin{proof}
\paragraph{Step 1: the full normal restriction is semisimple.}
First take the E-socle of \(\mathcal Q\) over \(\mathbf F_p\).
It is nonzero, since a finite-dimensional module has a simple
submodule.  Normality of E makes its socle H-invariant.  The
original-field irreducibility of H therefore makes that socle all
of \(\mathcal Q\).

We also need semisimplicity after scalar extension, and record a
direct descent argument to avoid any assumption about absolute
irreducibility.  Put \(F=\overline{\mathbf F}_p\) and
\(M=\mathcal Q\otimes F\).  The E-socle of M is nonzero,
H-invariant and stable under the semilinear p-power Frobenius:
the matrices of E and H have entries in \(\mathbf F_p\), and
Frobenius carries simple E-submodules to simple E-submodules.
A Frobenius-stable F-subspace descends to \(\mathbf F_p\).
Indeed its unique reduced-row-echelon basis is unchanged under
Frobenius, so all its coefficients lie in \(\mathbf F_p\).
The descended nonzero H-invariant subspace must be \(\mathcal Q\).
Thus M equals its E-socle and is semisimple.

There is no trivial E-constituent.  The original invariant space
\(\mathcal Q^E\) is H-invariant, so is zero or all of
\(\mathcal Q\).  The latter would contradict faithfulness on
the nontrivial E.  Solving the same linear invariant equations
after extension gives \(M^E=0\).

Every irreducible constituent S is therefore nontrivial.  Its
projective kernel in E is exactly Z(E).  To see this, the subgroup
of elements acting as scalars is normal in E and contains Z(E)
by Schur's lemma.  If its image in the simple \(A_n\) were
nontrivial, that image would be all of \(A_n\).  Then E would
be that subgroup times Z(E), and all of E would act by scalars.
A scalar representation of a perfect group is trivial, contrary
to the choice of S.  This proves projective faithfulness of
\(A_n\) on S.  Accepted D-TD32, applied to this actual
projective action, gives
\[
             \dim_F S>n/5,\qquad\text{hence}\qquad d>n/5.
 \tag{ANE.5}
\]
No modular irreducible-degree classification or multiplier table
is used in this step.

\paragraph{Step 2: every isotypic orbit contributes.}
Fix h and put \(x=\psi(h)\), \(k=k(h)\).  The case k=0 is
immediate, so suppose k>0.  Decompose M into its canonical
E-isotypic spaces.  The operator h permutes these spaces because
E is normal.  Each orbit of this permutation gives an h-invariant
direct summand.

For an orbit of length \(\ell\ge2\), let D be the dimension
of its entire summand.  All \(\ell\) component dimensions are
\(D/\ell\).  For any eigenvalue target \(\lambda\ne0\),
an eigenvector is determined by its component in the first space:
the equations in the successive spaces determine all the other
components.  The final equation can only reduce the dimension.
Consequently
\[
 \dim\ker(h-\lambda\mid M_{\rm orbit})\le D/\ell,
 \qquad
 \operatorname{rank}(h-\lambda\mid M_{\rm orbit})
       \ge D/2\ge\frac{Dk}{24n}.
 \tag{ANE.6}
\]
The target zero has zero kernel since h is invertible.  This
argument also works when p divides \(\ell\); it does not
diagonalize the permutation of the components.

It remains to treat a fixed isotypic space, written by evaluation
as \(S\otimes_F U\), with \(S\) irreducible for E and U its
entire multiplicity space.  Let \(\rho:E\to\operatorname{GL}(S)\)
be the actual constituent.  Since h fixes its isomorphism type,
there is an invertible intertwiner B satisfying
\[
 B\rho(z)B^{-1}=\rho(hzh^{-1})\quad(z\in E),\qquad
 h\big|_{S\otimes U}=B\otimes A
 \tag{ANE.7}
\]
for some invertible A on U.  Explicitly, on
\(U=\operatorname{Hom}_E(S,M)\), A sends f to \(h f B^{-1}\).
The displayed intertwiner identity makes this again an E-map.
This proves the full multiplicity factorization; no multiplicity
operator has been omitted.

\paragraph{Step 3: finite normalization and the actual odd extension.}
Let L be the finite order of h.  The operator \(B^L\) commutes
with \(\rho(E)\), so is scalar by Schur's lemma.  Over F choose
a scalar making the L-th power of B equal to one, and absorb the
reciprocal scalar in A.  This does not change geometric projective
rank.  B is now finite order and normalizes the finite group
\(\widetilde E=\rho(E)\).  Thus
\(G_S=\langle\widetilde E,B\rangle\) is a finite group: it is
a union of at most L cosets of \(\widetilde E\).

The group \(\widetilde E\) is a perfect central extension of
\(A_n\).  Its projective simple quotient is the same actual
\(A_n\) as in (ANE.1), by the kernel argument in Step 1.
Conjugation by B on this quotient is exactly conjugation by x.
Conjugation of \(S_n\) on \(A_n\) is faithful for \(n\ge5\):
a permutation centralizing \(A_n\) must fix the unique fixed
point of each natural point stabilizer, and hence fixes every
letter.  It follows that conjugation defines a surjection
\[
 G_S\longrightarrow\langle A_n,x\rangle,
 \qquad B\longmapsto x,
 \qquad\ker\text{ is scalar and central}.
 \tag{ANE.8}
\]
For the kernel assertion, an automorphism of the perfect
\(\widetilde E\) that induces the identity on
\(\widetilde E/Z(\widetilde E)\) differs from the identity by
a homomorphism into its center.  That homomorphism is trivial.
An element of the displayed kernel therefore centralizes
\(\widetilde E\); Schur's lemma makes its matrix scalar.

If x is even, choose \(z\in\widetilde E\) above x.  The
operators B and z induce the same automorphism of the perfect
\(\widetilde E\), by that same central-automorphism argument.
Their ratio is scalar.  D-TD31 applied inside \(\widetilde E\)
therefore bounds \(\nu_{\rm alg}(B)\).  If x is odd, (ANE.8) and the tautological
representation of \(G_S\) on S verify every hypothesis of
the precise odd-extension variant of D-TD31: the group is
finite, the kernel is central, its perfect core maps onto
\(A_n\), and restriction to that core is nontrivial and
irreducible.  In both cases the result is
\[
               \nu_{\rm alg}(B)\ge
                   \frac{(\dim S)k}{24n}.
 \tag{ANE.9}
\]
In particular a semilinear or Galois permutation of constituents
has not been assumed to preserve one constituent.  Nonfixed
isotypic spaces were already treated in Step 2.

\paragraph{Step 4: tensor rank and summation.}
The all-characteristic geometric tensor inequality in D-TD27
gives, for every scalar target,
\[
 \operatorname{rank}(B\otimes A-\lambda I)
      \ge (\dim U)\nu_{\rm alg}(B)
      \ge\frac{(\dim S)(\dim U)k}{24n}.
 \tag{ANE.10}
\]
It is applicable even when either operator has nontrivial Jordan
blocks.  One may view the kernel as an intertwiner space between
\(\lambda A^{-\mathsf t}\) and B; matching Jordan blocks of
sizes u,v contribute \(\min(u,v)\), which proves that bound.
Sum (ANE.6) and (ANE.10) over every isotypic orbit for the same
target \(\lambda\).  Their dimensions sum to d, giving the
projective rank bound in (ANE.3).  Taking \(\lambda=1\) and
using invariance of rank under field extension gives its ordinary
\(\mathbf F_p\)-rank consequence.

\paragraph{Step 5: the actual Heisenberg trace and cubic comparison.}
Assume now (ANE.4) and \(n\ge10^6\).  Compose the actual
surjections \(J\to P\to H\to T\), and call the result
\(x:J\twoheadrightarrow T\).  The accepted Weyl-basis trace
identity D-TD05 applies to every projective normalizer lift,
including its translation and scalar phases.  For a unitary
representative on W inducing h with \(k(h)>0\), it gives
\[
 \frac{|\operatorname{tr}_W A_h|}{a}
    \le p^{-\operatorname{rank}(1-h)/2}
    \le p^{-d k(h)/(48n)}<2^{-k(h)/240},
 \qquad
 \frac{|\operatorname{tr}_W A_h|}{a}\le2^{-1/2}<\frac78.
 \tag{ANE.11}
\]
The strict inequality uses \(d>n/5\); the second bound uses
faithfulness of the label action, since nontrivial natural image
implies \(h\ne1\) and hence \(\operatorname{rank}(1-h)\ge1\).

Write \(f=n-k\).  If \(0<k\le n/16\), then
\(2^{-1/2}<(15/16)^3\le(f/n)^3\); the first comparison is
equivalent to the exact integer inequality \(2\cdot15^6>16^6\).
If \(k>n/16\), then \(2^{-k/240}<2^{-n/3840}<n^{-3}\).
For the last strict inequality, choose the integer j with
\(2^j\le n<2^{j+1}\).  Here \(j\ge19\), and
\[
 2^{19}>11520\cdot20,\qquad
 2^j>11520(j+1)\ \Longrightarrow\
 2^{j+1}>11520(j+2).
 \tag{ANE.12}
\]
Therefore \(n/3840>3(j+1)\) and the claimed inequality follows.
The same estimate, with the larger exponent \(n/10\), gives
\(a=p^{d/2}>2^{n/10}>n^3\).

The entire external partner is retained in the trace comparison.
One can choose invariant Hermitian forms for the finite projective
factor images: a finite projective image has a finite preimage
in the special linear group, and averaging over it supplies such
a form.  For any actual coefficient operator write it as the
tensor of unitary representatives of those two projective images,
times a scalar.  Its finite order makes this scalar have absolute
value one.  The normalized trace of the whole operator is thus
at most that of the W-factor.  Apply this statement to every
original group element, and hence to every full original power.
For natural image one, the unitary trace bound one suffices.
Combining the two cases proves
\[
 \frac{|\operatorname{tr}_V(g)|}{e}
       \le\left(\frac{\max\{f(x(g)),1\}}n\right)^3,
 \qquad x(g)\ne1\ \Longrightarrow\quad
       \frac{|\operatorname{tr}_V(g)|}{e}<\frac78,
 \qquad e\ge a>n^3.
 \tag{ANE.13}
\]
These are exactly the two full-power whole-coefficient hypotheses
of accepted HD-FT27, with its actual alternating/symmetric quotient
and actual connected sphere cover.  That proposition contradicts
moving-rank sum \(2e+6\), proving the source conclusion.
\end{proof}

\paragraph{Retained existence boundary.}
The theorem begins with an actual normal quasisimple subgroup E
inside the actual label image H.  It does not construct E from
an abstract alternating quotient, and it does not assert that
the normal kernel is small.  D-TD29 supplies such an E in its
separately stated large alternating homogeneous-multiplicity
range.  Other kernel and symmetric sign cases still need their
own existence argument.  The result proves neither a geometric
source nor an original-character-field or polarization statement.
